% ============================================================================
% WORKBOOK — Section 7.2: Circumference of a Circle
% CCSS 7.G.B.4
% Practice-focused reinforcement for the study guide topic
% ============================================================================
\section{Circumference of a Circle}
\topicTitle{Circumference of a Circle}

\begin{quickReview}{Circumference of a Circle}
The \textbf{circumference} is the distance around a circle. Two formulas:

\begin{itemize}[leftmargin=*]
    \item $C = \pi d$ \quad (when you know the \textbf{diameter})
    \item $C = 2\pi r$ \quad (when you know the \textbf{radius})
\end{itemize}

\medskip
\textbf{Key fact:} $\pi \approx 3.14$. The ratio $\dfrac{C}{d}$ is always $\pi$ for every circle.

\smallskip
\textbf{Example:} A circle has radius $6$ cm. $C = 2 \times 3.14 \times 6 = 37.68$ cm.
\end{quickReview}

% ── Warm-Up ──────────────────────────────────────────────────────────────────
\begin{practiceBox}[funGreen]{\faSun~Warm-Up}

\practiceHeader[funGreen]{Find the Circumference (use $\pi \approx 3.14$)}

\begin{multicols}{2}
\prob Diameter $= 10$ cm \answerBlank[3cm]
\answerExplain{$31.4$ cm}{$C = \pi d = 3.14 \times 10 = 31.4$ cm.}

\prob Radius $= 4$ in \answerBlank[3cm]
\answerExplain{$25.12$ in}{$C = 2\pi r = 2 \times 3.14 \times 4 = 25.12$ in.}

\prob Diameter $= 6$ m \answerBlank[3cm]
\answerExplain{$18.84$ m}{$C = \pi d = 3.14 \times 6 = 18.84$ m.}

\prob Radius $= 10$ ft \answerBlank[3cm]
\answerExplain{$62.8$ ft}{$C = 2\pi r = 2 \times 3.14 \times 10 = 62.8$ ft.}

\prob Diameter $= 3$ cm \answerBlank[3cm]
\answerExplain{$9.42$ cm}{$C = \pi d = 3.14 \times 3 = 9.42$ cm.}

\prob Radius $= 7$ m \answerBlank[3cm]
\answerExplain{$43.96$ m}{$C = 2\pi r = 2 \times 3.14 \times 7 = 43.96$ m.}
\end{multicols}
\end{practiceBox}

% ── Practice A: Circumference from Diameter ──────────────────────────────────
\begin{practiceBox}{Circumference from Diameter}

\practiceHeader{Use $C = \pi d$ with $\pi \approx 3.14$. Show your work.}

\prob A dinner plate has a diameter of $25$ cm. Find its circumference. \answerBlank[3cm]
\answerExplain{$78.5$ cm}{$C = \pi d = 3.14 \times 25 = 78.5$ cm.}

\prob A circular mirror has a diameter of $40$ cm. What is its circumference? \answerBlank[3cm]
\answerExplain{$125.6$ cm}{$C = \pi d = 3.14 \times 40 = 125.6$ cm.}

\prob A tree trunk has a diameter of $1.5$ m. Find the circumference. \answerBlank[3cm]
\answerExplain{$4.71$ m}{$C = \pi d = 3.14 \times 1.5 = 4.71$ m.}

\prob A coin has a diameter of $2$ cm. What is the distance around it? \answerBlank[3cm]
\answerExplain{$6.28$ cm}{$C = \pi d = 3.14 \times 2 = 6.28$ cm.}

\prob A circular fountain has a diameter of $8.5$ m. Find the circumference. \answerBlank[3cm]
\answerExplain{$26.69$ m}{$C = \pi d = 3.14 \times 8.5 = 26.69$ m.}

\end{practiceBox}

% ── Practice B: Circumference from Radius ────────────────────────────────────
\begin{practiceBox}[funTeal]{Circumference from Radius}

\practiceHeader[funTeal]{Use $C = 2\pi r$ with $\pi \approx 3.14$.}

\prob $r = 5$ cm \answerBlank[3cm]
\answerExplain{$31.4$ cm}{$C = 2\pi r = 2 \times 3.14 \times 5 = 31.4$ cm.}

\prob $r = 12$ in \answerBlank[3cm]
\answerExplain{$75.36$ in}{$C = 2\pi r = 2 \times 3.14 \times 12 = 75.36$ in.}

\prob $r = 3.5$ m \answerBlank[3cm]
\answerExplain{$21.98$ m}{$C = 2\pi r = 2 \times 3.14 \times 3.5 = 21.98$ m.}

\prob $r = 20$ ft \answerBlank[3cm]
\answerExplain{$125.6$ ft}{$C = 2\pi r = 2 \times 3.14 \times 20 = 125.6$ ft.}

\prob $r = 0.5$ cm \answerBlank[3cm]
\answerExplain{$3.14$ cm}{$C = 2\pi r = 2 \times 3.14 \times 0.5 = 3.14$ cm.}

\end{practiceBox}

% ── Practice C: Working Backwards ────────────────────────────────────────────
\begin{practiceBox}[funOrange]{Find the Diameter or Radius}

\practiceHeader[funOrange]{Given the circumference, find the missing measurement. Use $\pi \approx 3.14$.}

\prob A circle has circumference $62.8$ cm. What is its diameter? \answerBlank[3cm]
\answerExplain{$20$ cm}{$d = C \div \pi = 62.8 \div 3.14 = 20$ cm.}

\prob A circle has circumference $18.84$ m. What is its radius? \answerBlank[3cm]
\answerExplain{$3$ m}{$d = C \div \pi = 18.84 \div 3.14 = 6$ m. Then $r = 6 \div 2 = 3$ m.}

\prob A circular garden border is $47.1$ ft long. What is the diameter of the garden? \answerBlank[3cm]
\answerExplain{$15$ ft}{$d = C \div \pi = 47.1 \div 3.14 = 15$ ft.}

\end{practiceBox}

% ── Word Problems ────────────────────────────────────────────────────────────
\begin{practiceBox}[funBlue]{\faGlobe~Word Problems}

\wordProblem{A bike wheel has a diameter of $70$ cm. How far does the bike travel in one full rotation of the wheel? Use $\pi \approx 3.14$.}{cm}
\answerExplain{$219.8$ cm}{One rotation covers the circumference: $C = \pi d = 3.14 \times 70 = 219.8$ cm.}

\wordProblem{A circular track has a radius of $50$ m. Mia jogs around the track $3$ times. How far does she run in total? Use $\pi \approx 3.14$.}{m}
\answerExplain{$942$ m}{One lap: $C = 2\pi r = 2 \times 3.14 \times 50 = 314$ m. Three laps: $314 \times 3 = 942$ m.}

\wordProblem{A gardener wants to put a fence around a circular flower bed with a diameter of $4$ m. Fencing costs $\$12$ per meter. How much will the fence cost? Use $\pi \approx 3.14$.}{dollars}
\answerExplain{$\$150.72$}{Circumference: $C = 3.14 \times 4 = 12.56$ m. Cost: $12.56 \times 12 = \$150.72$.}

\end{practiceBox}

% ── Challenge ────────────────────────────────────────────────────────────────
\begin{challengeBox}
\prob A semicircular arch has a diameter of $12$ m. What is the total length around the arch (the curved part plus the straight bottom)? Use $\pi \approx 3.14$. \answerBlank[3cm]
\answerExplain{$30.84$ m}{The curved part is half the circumference: $\frac{1}{2} \times 3.14 \times 12 = 18.84$ m. The straight bottom is the diameter: $12$ m. Total: $18.84 + 12 = 30.84$ m.}

\prob A clock has a minute hand that is $14$ cm long. How far does the tip of the minute hand travel in one full hour? Use $\pi \approx 3.14$. \answerBlank[3cm]
\answerExplain{$87.92$ cm}{The tip traces a circle with radius $14$ cm. $C = 2\pi r = 2 \times 3.14 \times 14 = 87.92$ cm.}
\end{challengeBox}

\encouragement{Fantastic work with circumference! You can now measure the distance around any circle!}
